\chapter{Conclusion}

The Deep Kalman Filter keeps the probabilistic structure of the Kalman
filter, a Markov chain of latent states with one emission per time step,
and makes the transition and emission non-linear with neural networks.
Exact inference is then impossible, and the model is trained like a VAE,
by maximising a lower bound on the log-likelihood.\\

Three points that the original paper leaves implicit are central to
understanding the method:
\begin{itemize}
    \item \textbf{The structure of the inference network follows from
    d-separation.} Given $\mathbf{z}_{t-1}$, the true posterior of
    $\mathbf{z}_t$ depends only on the future observations
    $\mathbf{x}_{t:T}$, which is why the inference network combines
    $\mathbf{z}_{t-1}$ with a backward RNN.
    \item \textbf{The ELBO splits into one KL per time step}, by the
    chain rule of the KL divergence. Each of these is between two
    Gaussians given $\mathbf{z}_{t-1}$ and has a closed form.
    \item \textbf{One sampled trajectory gives an unbiased estimate of
    the ELBO.} The ELBO is linear in its expectations under the
    intractable marginals of the inference network, and ancestral
    sampling provides exact samples of all of these marginals at once.
\end{itemize}

The autoregressive variants I designed add edges from past observations
to the latent state, which gives the transition model direct access to
the history and greatly improves forecasting on the sine-wave data.
Redoing the d-separation argument on the new graph shows that these
edges also change the posterior: it now depends on the past observations
as well as the future ones. The FAR DKF, whose inference network sees
both, is therefore the principled inference model for the AR DKF
generative model. On the Lorenz data, however, the three autoregressive
variants perform comparably: over four training runs per experiment,
neither the forecasting errors nor the ELBOs on test data rank them
consistently, and the variability between runs is larger than the
differences between models.\\

Several directions remain open. The variability between training runs
should be reduced, or at least controlled with fixed random seeds and
more repetitions, before the variants can be compared. Tighter estimates
of the log-likelihood, with many more importance samples, would be
needed to measure the differences between their ELBO gaps.
Forecasts should be evaluated with the predictive mean, or with a
probabilistic score, rather than with a single sampled path. The
reconstruction error currently matches the noise level, which suggests
that the inference network reproduces the observations, noise included,
rather than denoising them: monitoring the two terms of the ELBO
separately would show how much of the ELBO each part of the model
accounts for. Finally, control inputs, for example from an inertial
measurement unit, could be added to the transition model as in the
original formulation.
